\section{扩展 Whittaker 范畴}
\markboth{扩展 Whittaker 范畴}{}

现在转向本文的主角，扩展 Whittaker 范畴，记作 $\Whit(G)\ext_{glob}$。

Whittaker 范畴 $\Whit(G)\rd_{glob}$ 理应从 $D(\Bun_G)$ 接受一个函子，它「抓住」自守 D-模的全部非退化 Fourier 系数（在 $G$ 的中心连通的假设下）。扩展 Whittaker 范畴 $\Whit(G)\ext_{glob}$ 则应当抓住全部 Fourier 系数。

让我们先解释函数论上的类比。令 $\Char^{\mathrm{n.d.}}$ 记 $N(\A)$ 上在 $N(K)$ 上平凡的非退化特征所构成的空间。$T(\A)$ 在 $N(\A)$ 上的伴随作用给出 $T(K)/Z_G(K)$ 在 $\Char^{\mathrm{n.d.}}$ 上的作用，而这个作用是单可迁的（依据 $Z_G$ 连通的假设）。于是我们得到同构
\begin{equation}
  Z_G(K)N(K)\backslash G(\A)/G(\mathbb O)
  \simeq
  T(K)\backslash\Bigl(N(K)\backslash G(\A)/G(\mathbb O)\times\Char^{\mathrm{n.d.}}\Bigr).
  \label{eq:7.1}
\end{equation}

此外，在 $N(K)\backslash G(\A)/G(\mathbb O)\times\Char^{\mathrm{n.d.}}$ 上的函数空间上，可以施加自然的 $N(\A)$-等变条件（该条件依赖于 $\Char^{\mathrm{n.d.}}$ 的一个点），而这个条件在 $T(K)$-作用下不变。所得的函数空间因而可以等同于 $Z_G(K)N(K)\backslash G(\A)/G(\mathbb O)$ 上的 Whittaker 函数空间。

现在终于可以解释 $D(\Bun^{\mathrm{rat},\mathrm{red},\mathrm{ext}}_N)$ 及其满子范畴 $\Whit(G)\ext_{glob}$ 在函数论上的类比了。具体地说，令 $\Char$ 记 $N(\A)$ 上在 $N(K)$ 上平凡的全部特征所构成的空间。范畴 $D(\Bun^{\mathrm{rat},\mathrm{red},\mathrm{ext}}_N)$ 是
\[
  T(K)\backslash\Bigl(N(K)\backslash G(\A)/G(\mathbb O)\times\Char\Bigr)
\]
上函数空间的几何类比；而 $\Whit(G)\ext_{glob}$ 是其中「要求关于 $N(\A)$ 按依赖于 $\Char$ 的点的特征作等变」所得函数子空间的几何类比。

特别地要注意：当我们特殊到 $\{0\}\subset\Char$ 时，就回到
\[
  T(K)N(\A)\backslash G(\A)/G(\mathbb O)
\]
上的函数集合。

\begin{annote}
$\Char$ 与 $\Char^{\mathrm{n.d.}}$ 的区别是本节的关键：$\Char^{\mathrm{n.d.}}$ 只是非退化特征那一片，$T(K)$ 在其上单可迁地作用，故 (7.1) 把「$Z_G$ 处的 Whittaker 函数空间」等同于「$T(K)$ 作用下的非退化特征丛的截面」；而 $\Char$ 包含退化特征，其上的等变条件会退化，因而得到的是「全部 Fourier 系数」而非只有非退化部分。$D(\Bun^{\mathrm{rat},\mathrm{red},\mathrm{ext}}_N)$ 正是 $\Char$ 全体的几何对应物。\par\vspace{0.15em}{\raggedleft 【来源：原书 p.12 §7（A）；经典 Whittaker 模型中「非退化特征」与「全部特征」的对比（A）】\par}
\end{annote}

\ob{7.1}{blk:7.1}\textbf{范畴。}空间 $\overline{\Bun}^{\mathrm{red},\mathrm{ext}}_N$ 与范畴
\[
  D(\Bun^{\mathrm{rat},\mathrm{red},\mathrm{ext}}_N)
  := D\bigl(\overline{\Bun}^{\mathrm{red},\mathrm{ext}}_N\bigr)^{\mathcal H_{G,N}}
\]
的定义，与 $\overline{\Bun}\rd_N$ 和 $D(\Bun^{\mathrm{rat},\mathrm{red}}_N)^{\mathcal H_{G,N}}$ 的定义逐字相同，唯一的区别是：现在映射
\[
  \gamma_i^{\mathrm{red},\mathrm{ext}} : \Ch\alpha_i(P_T)\to\omega_X
\]
允许取零。

与上面一样，我们有伴随函子对
\[
  \ind_N^{\mathrm{red},\mathrm{ext}} :
  D\bigl(\overline{\Bun}^{\mathrm{red},\mathrm{ext}}_N\bigr)
  \rightleftarrows
  D(\Bun^{\mathrm{rat},\mathrm{red},\mathrm{ext}}_N) : \oblv_N^{\mathrm{red},\mathrm{ext}}.
\]

先前的讨论对群胚 $\mathcal H_N$ 同样适用，于是得到相应的满子范畴
\[
  \Whit(G)\ext_{glob} :=
  D\bigl(\overline{\Bun}^{\mathrm{red},\mathrm{ext}}_N\bigr)^{(\mathcal H_N,\chi),\mathcal H_{G,N}}
  \subset D\bigl(\overline{\Bun}^{\mathrm{red},\mathrm{ext}}_N\bigr)^{\mathcal H_{G,N}}
  =: D(\Bun^{\mathrm{rat},\mathrm{red},\mathrm{ext}}_N).
\]

我们有相应的伴随函子对
\[
  \ind_{\Whit\ext_{glob}} : D(\Bun^{\mathrm{rat},\mathrm{red},\mathrm{ext}}_N)
  \rightleftarrows \Whit(G)\ext_{glob} : \oblv_{\Whit\ext},
\]
其右伴随全忠实。

\ob{7.2}{blk:7.2}\textbf{扩展 Whittaker 范畴与 $D(\Bun\rt_T)$。}注意，把 $\gamma_i\rd$ 的资料置零，$\bBun_B$ 就闭嵌入到 $\overline{\Bun}^{\mathrm{red},\mathrm{ext}}_N$ 中。把这个映射记作 $\fI$。

我们把 $\fI$ 看作诱导出「闭嵌入」$\fI\rt : \Bun\rt_B\hookrightarrow \Bun^{\mathrm{rat},\mathrm{red},\mathrm{ext}}_N$，意思是说，我们有伴随函子对
\[
  (\fI\rt)_! := (\fI\rt)_* :
  D(\Bun\rt_B)\rightleftarrows D(\Bun^{\mathrm{rat},\mathrm{red},\mathrm{ext}}_N) : (\fI\rt)^!,
\]
其中 $(\fI\rt)_!$ 全忠实，并且使下图交换：
\[
\begin{tikzcd}[row sep=2.9em, column sep=3.4em, ampersand replacement=\&]
  D(\Bun\rt_B) \arrow[r] \arrow[d, "\oblv_B"'] \&
  D(\Bun^{\mathrm{rat},\mathrm{red},\mathrm{ext}}_N) \arrow[d, "\oblv_N^{\mathrm{red},\mathrm{ext}}"] \\
  D(\bBun_B) \arrow[r] \&
  D(\overline{\Bun}^{\mathrm{red},\mathrm{ext}}_N)
\end{tikzcd}
\]
（以及把相应函子换成其伴随后所得的另外三个图）。

上述伴随对把相应的 $\mathcal H_N$-等变范畴互相对应，从而诱导出伴随对
\[
  (\fI\rt)_! := (\fI\rt)_* :
  D(\Bun\rt_T)\rightleftarrows \Whit(G)\ext_{glob} : (\fI\rt)^!.
\]

\ob{7.3}{blk:7.3}\textbf{扩展与非扩展的范畴。}我们有开嵌入
\[
  \fJ : \overline{\Bun}\rd_N\hookrightarrow\overline{\Bun}^{\mathrm{red},\mathrm{ext}}_N .
\]
我们把 $\fJ$ 看作诱导出「开嵌入」
\[
  \fJ\rt : \Bun^{\mathrm{rat},\mathrm{red}}_N\hookrightarrow\Bun^{\mathrm{rat},\mathrm{red},\mathrm{ext}}_N,
\]
它对应的是所有 $\gamma_i\rd$ 均非零的那个轨迹。

严格地说，这意味着我们有如下限制函子
\[
  (\fJ\rt)^! := (\fJ\rt)^* :
  D(\Bun^{\mathrm{rat},\mathrm{red},\mathrm{ext}}_N)\to D(\Bun^{\mathrm{rat},\mathrm{red}}_N),
\]
它有右伴随，记作 $(\fJ\rt)_*$，且该右伴随全忠实。

这些函子就函子
\[
  \ind_{N^{\mathrm{red},\mathrm{ext}}} : D(\overline{\Bun}^{\mathrm{red},\mathrm{ext}}_N)
  \rightleftarrows D(\Bun^{\mathrm{rat},\mathrm{red},\mathrm{ext}}_N) : \oblv_{N^{\mathrm{red},\mathrm{ext}}}
\]
与
\[
  \ind_{N\rd} : D(\overline{\Bun}\rd_N)
  \rightleftarrows D(\Bun^{\mathrm{rat},\mathrm{red}}_N) : \oblv_{N\rd}
\]
而言，给出四个交换图。

\ob{7.3.1}{blk:7.3.1}上述函子对把相应的 Whittaker 子范畴互相对应，于是我们得到伴随对
\[
  (\fJ\rt)^* : \Whit(G)\rd_{glob}\rightleftarrows\Whit(G)\ext_{glob} : (\fJ\rt)_*,
\]
其中 $(\fJ\rt)_*$ 全忠实。

下面我们来证明：

\begin{othm}{命题}{7.3.2}{prop:7.3.2}
$(\fJ\rt)^!$ 的部分定义的左伴随（记作 $(\fJ\rt)_!$）在子范畴 $\Whit(G)\rd_{glob}\subset D(\Bun^{\mathrm{rat},\mathrm{red}}_N)$ 上有定义。
\end{othm}

由命题形式地推出：$(\fJ\rt)_!$ 把 $\Whit(G)\ext_{glob}$ 送到 $\Whit(G)\rd_{glob}$，并且全忠实，从而是 $(\fJ\rt)^!$ 的右逆。

\begin{proof}
把 Hecke 范畴看作 $\Rep(\Ch G,\mathrm{Ran}_X,un)$ 的 Ran 版本，它作用在所有涉及的范畴上。这个作用与函子 $(\fJ\rt)^!$ 交换；又因 $\Rep(\Ch G,\mathrm{Ran}_X,un)$ 是刚性的，它也与 $(\fJ\rt)_!$ 在后者有定义的子范畴上交换。

由定理 5.3.3 推出，$\Whit(G)\rd_{glob}$ 在 Hecke 作用下由单个对象生成，即由 $\one_{\Whit(G)\rd_{glob}}$ 生成。因此只需证明
\[
  (\fJ\rt)_!\bigl(\one_{\Whit(G)\rd_{glob}}\bigr)
\]
有定义。而后者显然如此：$\one_{\Whit(G)\rd_{glob}}$ 是把 $\ind_N^{\mathrm{red},\mathrm{ext}}$ 作用于 $D(\overline{\Bun}^{\mathrm{red},\mathrm{ext}}_N)$ 中一个完整（holonomic）对象而得。
\end{proof}

\begin{othm}{拟命题}{7.3.3}{prop:7.3.3}
函子 $(\fI\rt)_* : D(\Bun\rt_T)\to\Whit(G)\ext_{glob}$ 有左伴随。
\end{othm}

我们只在 $G$ 的半单秩为 $1$ 时证明本命题：

\begin{proof}
对对象 $\mathcal F\in\Whit(G)\ext_{glob}$，考虑好三角
\[
  (\fJ\rt)_!\circ(\fJ\rt)^!(\mathcal F)\to\mathcal F\to\mathcal F_1,
\]
其中 $\mathcal F_1$ 满足 $(\fJ\rt)^!(\mathcal F_1)=0$。

然而，当 $G$ 的半单秩为 $1$ 时，后一条件意味着 $\mathcal F_1$ 落在 $(\fI\rt)_*$ 的本质像中，于是函子 $(\fI\rt)^*$ 在它上面有定义。
\end{proof}

\begin{annote}
两处省略需要点破。7.3.2 的证明里，「在 Hecke 作用下由单个对象生成」用的是定理 5.3.3 把 $\Whit(G)\rd_{glob}$ 与 $\Rep(\Ch G,\mathrm{Ran}_X,un)$ 等同起来，而后者作为刚性幺半范畴在其正则对象下生成自身；「刚性」是为了让左伴随与 Hecke 作用交换。7.3.3 的证明里，半单秩 $1$ 只有两个分层（$\fI$ 的像与 $\fJ$ 的像），故 $\mathcal F_1$ 必落在 $\fI$ 的像上，$(\fI\rt)^*$ 才有定义。\par\vspace{0.15em}{\raggedleft 【来源：原书 p.14 §7.3.2--7.3.3（A）；译者按（用两个分层的情形补足原文省略的推理）】\par}
\end{annote}

\ob{7.4}{blk:7.4}\textbf{系数函子。}我们有自然的遗忘映射
\[
  \fR : \overline{\Bun}^{\mathrm{red},\mathrm{ext}}_N\to\bBun_B .
\]
拉回诱导函子
\[
  \fR^! : D(\bBun_B)\to D(\overline{\Bun}^{\mathrm{red},\mathrm{ext}}_N).
\]
我们把 $\fR$ 看作诱导出「叠」的映射 $\Bun^{\mathrm{rat},\mathrm{red},\mathrm{ext}}_N\to\Bun\rt_B$。严格地说，我们的意思是存在典范定义的函子
\[
  (\fR\rt)^! : D(\Bun\rt_B)\to D(\Bun^{\mathrm{rat},\mathrm{red},\mathrm{ext}}_N),
\]
它使下图交换：
\[
\begin{tikzcd}[row sep=2.9em, column sep=4.2em, ampersand replacement=\&]
  D(\Bun\rt_B) \arrow[r, "(\fR\rt)^!"] \arrow[d, "\oblv_B"'] \&
  D(\Bun^{\mathrm{rat},\mathrm{red},\mathrm{ext}}_N) \arrow[d, "\oblv_N^{\mathrm{red},\mathrm{ext}}"] \\
  D(\bBun_B) \arrow[r, "\fR^!"'] \&
  D(\overline{\Bun}^{\mathrm{red},\mathrm{ext}}_N)
\end{tikzcd}
\]
并且把纵向箭头分别换成 $\ind_B$ 与 $\ind_N^{\mathrm{red},\mathrm{ext}}$ 后的另一个图也交换。

我们定义函子
\[
  \coeff\ext : D(\Bun_G)\to\Whit(G)\ext_{glob}
\]
为复合
\[
  \coeff\ext := \ind_{\Whit\ext,glob}\circ(\fR\rt)^!\circ(\fP\rt)^! .
\]

复合
\[
  \coeff\rd := (\fJ\rt)^!\circ\coeff\ext : D(\Bun_G)\to\Whit(G)\rd_{glob}
\]
是通常的 Whittaker 系数函子。

注意，复合
\[
  (\fI\rt)^!\circ\coeff\ext : D(\Bun_G)\to D(\Bun\rt_T)
\]
与函子 $\CT\rt$ 一致。

\begin{othm}{拟命题}{7.4.1}{prop:7.4.1}
函子 $\coeff\ext$ 有左伴随。
\end{othm}

我们把这一左伴随函子记作 $\Poinc\ext_!$。复合
\[
  \Poinc\rd_! := \Poinc\ext_!\circ(\fJ\rt)_! :
  \Whit(G)\rd_{glob}\to D(\Bun_G)
\]
是 $\coeff\rd$ 的左伴随，即通常的 Poincaré 级数函子。

下面我们只在 $G$ 的半单秩为 $1$ 时证明本定理。

\begin{proof}
如同定理 7.3.3 的证明，只需分别证明 $\Poinc\rd_!$ 有定义，以及 $\CT\rt$ 的左伴随有定义。后者已在 §3.2 中讨论过。前者如命题 7.3.2 的证明。
\end{proof}

\begin{annote}
原文此处称「定理 7.3.3」，但 7.3.3 是拟命题，7.3.2 才是命题；译文按原样译出并在工程文件疑点清单中记录。\par\vspace{0.15em}{\raggedleft 【来源：原书 p.15 §7.4（A）】\par}
\end{annote}

\ob{7.5}{blk:7.5}\textbf{全忠实性。}我们处理 $GL_n$ 情形下几何朗兰兹的方法，关键在于下述论断：

\begin{othm}{拟定理}{7.5.1}{thm:7.5.1}
当 $G=GL_n$ 时，函子 $\coeff\ext$ 全忠实。
\end{othm}

对其他的约化群，我们不知道类似的论断是否成立。

\ob{7.5.2}{blk:7.5.2}下面朝定理 7.5.1 的证明走一步。

对约化群 $G$ 与抛物子群 $P$，可以仿照 $D(\Bun\rt_B)$（用部分的 Plücker 数据）定义范畴 $D(\Bun\rt_P)$。我们有伴随函子对
\[
  (\fP\rt_P)_! : D(\Bun\rt_P)\rightleftarrows D(\Bun_G) : (\fP\rt_P)^!,
\]
并且如同定理 1.3.1，我们有：

\begin{othm}{拟定理}{7.5.3}{thm:7.5.3}
函子 $(\fP\rt_P)^!$ 全忠实。
\end{othm}

现在取 $G=GL_n$，$P$ 为抛物子群 $(n-1,1)$。下述论断「应当」可以用经典的迭代 Fourier 变换手续证明：

\begin{othm}{拟定理}{7.5.4}{thm:7.5.4}
存在范畴的等价
\[
  D(\Bun\rt_P)\simeq\Whit(G)\ext_{glob}.
\]
\end{othm}

复合
\[
  D(\Bun_G)\xrightarrow{(\fP\rt_P)^!} D(\Bun\rt_P)\to\Whit(G)\ext_{glob}
\]
就是函子 $\coeff\ext$。

显然，上述两个定理蕴含定理 7.5.1。

\begin{annote}
这三步的用意是：把「$G$ 的扩展 Whittaker 范畴」换成「抛物 $P$ 的有理范畴」，而后者是更接近 $D(\Bun_G)$ 的对象。拟定理 7.5.3 说 $(\fP\rt_P)^!$ 全忠实，拟定理 7.5.4 说 $D(\Bun\rt_P)$ 与扩展 Whittaker 范畴等价；两者一夹，就得到拟定理 7.5.1。原文把 7.5.1、7.5.3、7.5.4 都标为拟定理，后文却简称「定理」，译文照录其简称。\par\vspace{0.15em}{\raggedleft 【来源：原书 p.16 §7.5.2（A）；译者按（把三步的逻辑关系理清）】\par}
\end{annote}
